\section{Systolic Array Architectures}

\begin{remark}
    Systolic arrays might be hard to understand intuitively. The key idea is to have a \textbf{very specific architecture} of processing elements (PEs) that are connected in a \textbf{grid- or array-like structure} and control the flow of data. This allows us to process large amounts of data in parallel without having to load data entries multiple times.
\end{remark}

\begin{important}
    Systolic arrays are \textbf{not a standard pipeline}, they allow for \textbf{non linear and multi-dimensional data flow}.
\end{important}

\begin{example}
    To convolve a matrix with a kernel, we can use a systolic array to process the data in parallel. 
\end{example}

\begin{remark}
    The main advantage of systolic arrays is their high efficiency in terms of data movement and processing power as well as the regularity of the design. 
\end{remark}

\begin{remark}
    The main disadvantage is that they are very specialized and not good at exploiting irregular parallelism.
\end{remark}